% Discussion. Qualitative claims here must match data/agentic/RESULTS.md;
% every number is a macro from generated_numbers.tex.

\textbf{Measure self-inflicted harm, not only containment.} Containment is
a reasonable proxy on average, and that is what makes it dangerous: it is
right often enough to be trusted, and wrong on exactly the defenders whose
containment comes from stranding services. A containment leaderboard would
have ranked the containment-trained learner \numQContainmentTodayRankContainment\
of \numRankAgents; the hospital would rank it \numQContainmentTodayRankClinical.
Evaluations of autonomous defenders for critical services should report
the mission outcome alongside containment. They should also report the
self-inflicted share of that outcome, which only a dependency model can
compute, because it is the part of the harm the defender chose.

\textbf{Who may take the bet.} Our central result is that self-inflicted
outage is a gamble whose sign depends on the estate. An early cut that
strands services buys containment, and where detection is slow that
containment is worth far more than the stranded service. Where detection is
fast it buys almost nothing. A learner trained across estates takes the same
gamble everywhere. The dependency-closure shield forbids the gamble, which
protects every fast-detection hospital we simulated and forfeits large gains
in slow ones. Neither ``always'' nor ``never'' is right. The structural
answer is to make the shield the default and let a person who knows the
estate lift it: per deployment (``this hospital's detection is slow; the
defender may self-inflict''), or per incident, with the shield escalating a
blocked cut to an on-call operator rather than silently refusing it. The
agent proposes the gamble; a human with accountability for the estate
authorises it. The shield is what makes that allocation enforceable rather
than advisory. On these data, a shield enabled only in the fast-detection
profile beats both ``always'' and ``never'' for \numPosthocFastOnlyBestCount\
of the \numPosthocDefenders\ defenders. For the other
\numPosthocAlwaysBestCount, the disconnect playbook and the containment
learner on today's estate, ``always'' is better, because their cuts do harm
even in intermediate hospitals. So no fixed rule is right for every
defender, which is itself the case for authorising per deployment rather
than fixing a policy. We chose these rules after seeing the data and
evaluated them on the same data, so they motivate the design and are not
evidence for it.

\textbf{Hard constraint or soft penalty.} The soft-penalty learners show
what an availability term in the reward does: it prices self-inflicted
outage without preventing it (\numQSoftTodaySelfPct\ and
\numQSoftReplicasSelfPct\ of episodes). That can be the right behaviour,
and those learners are among the best on average. But a penalty offers no
guarantee for any particular hospital, and its weight is set before
deployment by someone who cannot know every estate it will run on. A shield
offers the guarantee and none of the judgement. The combination we would
deploy trains with shaping, so that good cuts are likely, and runs with an
operator-liftable shield, so that the judgement stays with a person.

\textbf{A dependency specification is a deployment prerequisite.} The
shield needs nothing but the service-dependency graph and the availability
rule. Those are exactly what an organisation must write down before any
automated responder may cut its network, and exactly what most hospitals do
not have in machine-readable form. Every intervention is an auditable event:
the responder asked to cut, and the dependency graph said the cut would
strand these named services. A wrong or incomplete specification is the
shield's failure mode, as the micro lockdown shows, so it must be versioned
and reviewed like any other safety case.

\textbf{Agent-agnostic by construction.} The shield screens proposed
response modes and never inspects the policy that proposed them. The same
monitor applies unchanged to a language-model responder, at the tool
boundary where the agent's ``disconnect'' call becomes a firewall change.
A language-model responder's reasoning cannot be verified; the effect of its
proposed action on a declared dependency graph can.

\textbf{Design the estate so the safe action is also the effective one.}
The shield removes the harmful cut; it cannot supply a harmless one. The
islands playbook cuts as hard as the disconnect playbook and loses
\numKeyDisconnectToIslandsHoursAbs\ fewer clinical service-hours, and on the
replica estate every shielded learner beats passive in every profile. The
cost of safe autonomy is partly an architectural property, decided long
before any incident.

\textbf{The shield is myopic, by design.} It forbids an action whose
\emph{immediate} effect removes service, even when accepting that loss would
pay off later. That is the source of both its guarantee and its price. A
shield that looked ahead would need a model of the attacker's future, and
its guarantee would inherit that model's errors. We chose checkability and
report what it costs; the operator override above is how the cost is
recovered where it is worth recovering.

\textbf{Beyond hospitals.} Any estate whose availability flows through
shared dependencies has the same structure: industrial control with a
central historian, a bank with one identity provider, an airline with one
reservation system. Wherever a containment cut can strand a service from
what it depends on, an autonomous defender needs the same guard and the
same question answered: who may accept the outage.
